\documentclass[11pt,a4paper]{ctexart}
\usepackage[a4paper,top=18mm,bottom=20mm,left=19mm,right=19mm]{geometry}
\usepackage{booktabs,tabularx,array,multirow}
\usepackage{xcolor,tcolorbox}
\usepackage{amsmath,amssymb}
\usepackage{graphicx}
\usepackage{pgfplots}
\usepackage{fancyhdr}
\usepackage{hyperref}
\usepackage{enumitem}
\usepackage{microtype}
\pgfplotsset{compat=1.18}

\definecolor{navy}{HTML}{17324D}
\definecolor{blue}{HTML}{2F6BFF}
\definecolor{teal}{HTML}{00A6A6}
\definecolor{orange}{HTML}{E9852D}
\definecolor{purple}{HTML}{7456B8}
\definecolor{rose}{HTML}{D45576}
\definecolor{green}{HTML}{3B9A66}
\definecolor{ink}{HTML}{263442}
\definecolor{muted}{HTML}{66788A}
\definecolor{panel}{HTML}{F3F6FA}
\definecolor{rule}{HTML}{D8E0E8}

\hypersetup{colorlinks=true,linkcolor=navy,urlcolor=blue,pdftitle={Sync2Act Quality-Aware ACT Mixed-Quality Ablation Study}}
\setlength{\parindent}{2em}
\setlength{\parskip}{0.25em}
\setlength{\headheight}{24pt}
\setlist{nosep,leftmargin=2em}
\renewcommand{\arraystretch}{1.22}
\pagestyle{fancy}
\fancyhf{}
\lhead{\textcolor{muted}{Sync2Act / Quality-Aware ACT}}
\rhead{\textcolor{muted}{v1.2.0 实验报告}}
\cfoot{\textcolor{muted}{\thepage}}
\renewcommand{\headrulewidth}{0.4pt}
\renewcommand{\headrule}{\hbox to\headwidth{\color{rule}\leaders\hrule height \headrulewidth\hfill}}

\newcolumntype{Y}{>{\raggedright\arraybackslash}X}
\newcolumntype{C}[1]{>{\centering\arraybackslash}p{#1}}
\newcommand{\code}[1]{\texttt{\small #1}}
\newcommand{\good}[1]{\textcolor{green}{\bfseries #1}}
\newcommand{\warn}[1]{\textcolor{orange}{\bfseries #1}}

\begin{document}

\begin{titlepage}
\pagecolor{panel}
\vspace*{15mm}
{\color{blue}\rule{25mm}{2.5pt}}\\[11mm]
{\fontsize{29}{35}\selectfont\bfseries\color{navy} Sync2Act\\[2mm]
Quality-Aware ACT\\[1mm]混合质量消融研究\par}
\vspace{5mm}
{\Large\color{muted}完整数据集、六组消融、三个随机种子\par}
\vspace{16mm}

\begin{tcolorbox}[colback=white,colframe=rule,boxrule=0.7pt,arc=2mm,left=5mm,right=5mm,top=4mm,bottom=4mm]
\begin{tabularx}{\linewidth}{>{\bfseries\color{navy}}p{33mm}Y}
实验规模 & 3 个真实单任务数据集 $\times$ 6 个模型消融 $\times$ 5 个训练条件 $\times$ 3 个随机种子\\
完成状态 & \good{270 / 270 组成功完成}\\
训练约束 & 完整 episode 与完整 frame；仅训练集损坏；验证集和测试集保持干净\\
核心指标 & 干净测试集 action MSE，以及相对同模型干净训练的 MSE 比值\\
生成日期 & 2026-09-13
\end{tabularx}
\end{tcolorbox}

\vfill
{\large\color{ink}\bfseries 结论预览}\par
\vspace{2mm}
{\color{ink}Oracle 模态质量对 temporal state shift 有稳定信号，但尚未形成跨损坏类型一致、幅度很大的收益。Quality-Full 相对 ACT-Lite 在 36 个“数据集 $\times$ 损坏条件 $\times$ seed”配对中有 22 个取得更低绝对 MSE；这说明方向值得继续验证，但不能据此声称闭环控制性能已提升。\par}
\vspace{12mm}
{\small\color{muted}实验输出：\code{runs/quality\_ablation\_v1\_2\_0} \hfill Sync2Act}
\end{titlepage}
\nopagecolor

\section{研究目标与结论边界}

本轮实验研究的不是“模型能否识别坏数据”，而是：当同一训练集中同时存在干净与损坏片段时，把按模态拆分的 Oracle quality 提供给 ACT-Lite，能否减弱坏示范对干净测试集动作预测的影响。Oracle quality 由已知的合成损坏参数直接计算，因此代表质量信息的理论上界，而不是部署时可直接获得的传感器输出。

所有结果均为\textbf{离线动作预测}。报告中的 MSE、MAE 和鲁棒性比值不能替代真实机器人或仿真器中的闭环成功率、回报和完成时间。因而本报告可以支持“训练数据损坏如何影响预测误差”的结论，但不能直接推出“机器人任务成功率提高”。

\section{数据集与训练协议}

三个数据集均为真实环境、单任务数据集；训练使用所有可用相机。数据按 episode 划分，避免同一 episode 的相邻帧同时进入训练与测试。

\begin{table}[htbp]
\centering
\small
\begin{tabularx}{\linewidth}{Y C{14mm} C{18mm} C{14mm} C{22mm} C{24mm}}
\toprule
数据集 & Episode & Frame & 相机 & State / Action & Train / Val / Test\\
\midrule
XArm Lift Medium & 800 & 20,000 & 1 & 4 / 4 & 640 / 80 / 80\\
Koch Pick Place 1 Lego & 100 & 32,951 & 2 & 6 / 6 & 80 / 10 / 10\\
ALOHA Static Battery & 49 & 29,400 & 4 & 14 / 14 & 41 / 4 / 4\\
\bottomrule
\end{tabularx}
\caption{完整数据集统计。图像输入分别来自 1、2、4 个相机。}
\end{table}

统一训练配置为 10 epochs、CUDA、随机种子 7/17/27、划分种子 7；ACT batch size 为 256，action segment length 为 16。图像最长边缩放至 64 像素。Normalization statistics 只从\textbf{干净训练集}计算并固定，避免损坏条件改变归一化标尺。

\begin{tcolorbox}[colback=panel,colframe=rule,title={\color{navy}\bfseries 数据隔离规则},fonttitle=\bfseries]
原始数据按 episode 划分后，仅 Training episodes 根据实验条件构造混合损坏；Validation episodes 与 Test episodes 始终保持干净。测试指标按 episode 内部计算轨迹量，再汇总，避免跨 Episode 跳变。
\end{tcolorbox}

\section{模态质量与混合质量构造}

每个样本不再共享单一 quality，而是包含 \code{image\_quality}、\code{state\_quality} 与 \code{action\_label\_quality}。同时分别记录模态级 \code{time\_offset}、\code{missing\_mask} 和 provenance，使模型可以区分“观测不可信”和“动作标签不可信”。时间错位的质量随严重度连续衰减：
\[
q_m=\exp\!\left(-\frac{|\Delta t_m|}{2\,\widetilde{\Delta t}}\right),
\]
其中 $\Delta t_m$ 为模态 $m$ 的时间偏移，$\widetilde{\Delta t}$ 为数据集的中位帧间隔。因此 3 秒 shift 的 quality 低于 1 秒 shift，而非二者都使用同一个常数。

单一损坏条件按连续片段构造 50\% 干净、50\% 损坏样本；混合条件包含 25\% 干净、25\% temporal state shift、25\% temporal action shift、25\% action delay。连续片段比逐帧独立采样更接近真实传感器或同步故障。

\begin{table}[htbp]
\centering
\small
\begin{tabularx}{\linewidth}{p{39mm}Y p{29mm}}
\toprule
训练条件 & 训练集构造 & 主要受影响质量\\
\midrule
Clean & 100\% 干净 & 全部为 1\\
Temporal state shift 2 & 50\% 干净 + 50\% state 时间错位 & state quality\\
Temporal action shift 2 & 50\% 干净 + 50\% action 时间错位 & action-label quality\\
Action delay 2 & 50\% 干净 + 50\% action 延迟 & action-label quality\\
Mixed three corruptions & 四类各占 25\% & 按损坏模态分别变化\\
\bottomrule
\end{tabularx}
\end{table}

\section{六组 Quality-Aware 消融}

\begin{table}[htbp]
\centering
\small
\begin{tabularx}{\linewidth}{p{38mm} C{23mm} C{25mm} Y}
\toprule
模型 & Quality 输入 & Quality 加权 loss & 目的\\
\midrule
ACT-Lite & 否 & 否 & 无质量信息的基线\\
Quality-Input & 是 & 否 & 测试质量作为条件输入是否有效\\
Quality-Weighted-Loss & 否 & 是 & 测试降低低质量标签权重是否有效\\
Quality-Full & 是 & 是 & 完整方法\\
Quality-Shuffled & 打乱 & 是 & 负对照：破坏质量与样本的对应关系\\
Quality-Constant & 常数 & 是 & 负对照：只有固定质量值\\
\bottomrule
\end{tabularx}
\end{table}

Quality-Shuffled 和 Quality-Constant 用于辨别收益究竟来自“真实的逐样本质量信号”，还是仅来自额外输入维度、损失缩放或随机训练差异。

\section{指标}

主要指标为干净测试集上的 action mean squared error：
\[
\mathrm{MSE}=\frac{1}{N d}\sum_{i=1}^{N}\lVert\hat{a}_i-a_i\rVert_2^2.
\]
由于三个数据集的动作单位和数值尺度不同，跨数据集汇总采用同一模型相对其干净训练条件的比值
\[
R=\frac{\mathrm{MSE}_{\mathrm{corrupt\ train}\rightarrow\mathrm{clean\ test}}}
{\mathrm{MSE}_{\mathrm{clean\ train}\rightarrow\mathrm{clean\ test}}}.
\]
$R=1$ 表示与干净训练持平；$R>1$ 表示坏示范使预测变差；$R<1$ 表示该训练条件下测试误差反而降低。比值先在每个数据集和 seed 内计算，再求均值。

\section{总体结果}

\begin{figure}[htbp]
\centering
\begin{tikzpicture}
\begin{axis}[
width=0.98\linewidth,height=72mm,
ymin=0.94,ymax=1.31,
ylabel={平均 MSE 比值 $R$},
symbolic x coords={Clean,State shift,Action shift,Action delay,Mixed},
xtick=data,x tick label style={rotate=18,anchor=east,font=\small},
legend style={at={(0.5,-0.31)},anchor=north,legend columns=3,font=\scriptsize,draw=none},
grid=major,grid style={rule!60},axis line style={rule},tick style={rule},
every axis plot/.append style={thick,mark size=2pt}]
\addplot[color=navy,mark=*] coordinates {(Clean,1.000000)(State shift,1.242326)(Action shift,1.004153)(Action delay,0.992422)(Mixed,1.008763)};
\addlegendentry{ACT-Lite}
\addplot[color=blue,mark=square*] coordinates {(Clean,1.000000)(State shift,1.154973)(Action shift,0.979456)(Action delay,0.967679)(Mixed,1.028503)};
\addlegendentry{Quality-Input}
\addplot[color=orange,mark=triangle*] coordinates {(Clean,1.000000)(State shift,1.240152)(Action shift,0.981021)(Action delay,0.981284)(Mixed,1.070521)};
\addlegendentry{Weighted-Loss}
\addplot[color=teal,mark=diamond*] coordinates {(Clean,1.000000)(State shift,1.159270)(Action shift,0.977767)(Action delay,0.981990)(Mixed,1.052382)};
\addlegendentry{Quality-Full}
\addplot[color=rose,mark=pentagon*] coordinates {(Clean,1.000000)(State shift,1.279684)(Action shift,0.990814)(Action delay,0.978826)(Mixed,1.026809)};
\addlegendentry{Shuffled}
\addplot[color=purple,mark=x] coordinates {(Clean,1.000000)(State shift,1.214256)(Action shift,0.978109)(Action delay,0.969221)(Mixed,0.997339)};
\addlegendentry{Constant}
\addplot[color=muted,dashed,no marks] coordinates {(Clean,1)(Mixed,1)};
\end{axis}
\end{tikzpicture}
\caption{三个数据集与三个 seeds 的平均相对 MSE。越低越好；虚线为干净训练基准。}
\end{figure}

\begin{table}[htbp]
\centering
\small
\begin{tabular}{lrrrrr}
\toprule
模型 & Clean & State shift & Action shift & Action delay & Mixed\\
\midrule
ACT-Lite & 1.000 & 1.242 & 1.004 & 0.992 & 1.009\\
Quality-Input & 1.000 & \good{1.155} & 0.979 & \good{0.968} & 1.029\\
Weighted-Loss & 1.000 & 1.240 & 0.981 & 0.981 & 1.071\\
Quality-Full & 1.000 & 1.159 & \good{0.978} & 0.982 & 1.052\\
Shuffled & 1.000 & 1.280 & 0.991 & 0.979 & 1.027\\
Constant & 1.000 & 1.214 & 0.978 & 0.969 & \good{0.997}\\
\bottomrule
\end{tabular}
\caption{平均 MSE 比值，列内最小值以绿色标出。}
\end{table}

Temporal state shift 是唯一在所有模型上都明显恶化的损坏类型，ACT-Lite 平均上升 24.2\%。Quality-Input 将该增幅降至 15.5\%，Quality-Full 为 15.9\%，说明显式质量输入对 state 错位最有希望。另一方面，Action shift、Action delay 和 Mixed 的均值大多接近 1；这些条件下约 1--3\% 的变化不足以构成稳健的总体收益。

\section{配对比较与数据集差异}

在 36 个损坏配对（3 数据集 $\times$ 4 损坏条件 $\times$ 3 seeds）中，Quality-Full 相对 ACT-Lite 有 22 个取得更低绝对 MSE，按归一化鲁棒性比较则为 23/36。Quality-Full 相对 Shuffled 为 25/36，相对 Constant 为 21/36。这些负对照说明正确对应的 quality 确实携带信号，但优势尚不够一致。

\begin{table}[htbp]
\centering
\small
\begin{tabularx}{\linewidth}{Y C{27mm} C{27mm} C{31mm}}
\toprule
比较 & 更低绝对 MSE & 更好归一化比值 & 总配对数\\
\midrule
Quality-Full vs ACT-Lite & 22 & 23 & 36\\
Quality-Full vs Shuffled & 25 & -- & 36\\
Quality-Full vs Constant & 21 & -- & 36\\
\bottomrule
\end{tabularx}
\end{table}

\begin{table}[htbp]
\centering
\small
\begin{tabularx}{\linewidth}{Y lrrrr}
\toprule
数据集 & 模型 & State shift & Action shift & Action delay & Mixed\\
\midrule
XArm Lift & ACT-Lite & 1.02 & 1.09 & 1.08 & 1.12\\
 & Quality-Full & 1.02 & 1.03 & 1.04 & 1.08\\
Koch Pick Place & ACT-Lite & 1.36 & 1.02 & 1.01 & 1.01\\
 & Quality-Full & 1.22 & 0.96 & 0.96 & 1.05\\
ALOHA Battery & ACT-Lite & 1.34 & 0.90 & 0.88 & 0.90\\
 & Quality-Full & 1.25 & 0.95 & 0.95 & 1.03\\
\bottomrule
\end{tabularx}
\caption{按数据集汇总的 MSE 比值。数值为报告中跨三个 seeds 的均值，保留两位小数。}
\end{table}

数据集差异很大：Temporal state shift 对 Koch 与 ALOHA 的影响最强，而 XArm 变化较小；ALOHA 在 action shift 与 delay 下出现低于 1 的比值，但这种“改善”没有在 XArm 和 Koch 上复现，可能来自损坏带来的正则化、动作冗余或随机波动，不应直接解释为延迟有益。

\section{结论、限制与下一步}

\begin{tcolorbox}[colback=panel,colframe=teal,boxrule=0.8pt,arc=1.5mm,title={\color{navy}\bfseries 主要结论}]
\begin{enumerate}
\item 数据损坏对训练的影响高度依赖损坏模态和数据集；Temporal state shift 是当前最稳定的负面因素。
\item Oracle 模态质量包含可利用信息。Quality-Input 和 Quality-Full 明显缓解了平均 state-shift 退化，且 Quality-Full 多数配对优于打乱质量的负对照。
\item Quality-Full 并未在所有损坏条件和数据集上一致胜出。因此当前证据支持“继续研究”，不支持“已解决坏示范训练问题”。
\end{enumerate}
\end{tcolorbox}

本研究仍有四项关键限制：（1）仅测试三个 seeds；（2）评价为离线动作误差，没有闭环 rollout；（3）跨数据集动作单位不同，绝对 MSE 不可直接横向比较；（4）quality 来自合成损坏真值，实际部署时未知。图像缩放到较低分辨率也可能削弱视觉质量差异。

建议下一阶段按以下顺序推进：
\begin{enumerate}
\item 先围绕 temporal state shift 增加严重度梯度与随机种子，检验收益是否稳定且随 quality 单调变化；
\item 将质量输入、加权函数及其校准误差进一步拆开，报告 bootstrap 置信区间和配对显著性；
\item 使用已知 corruption provenance 训练自动质量估计器，再用其预测质量替代 Oracle quality，衡量从上界到可部署方案的性能损失；
\item 最后在仿真或真实机器人上进行闭环 rollout，报告成功率与安全失败类型，再决定是否扩展至完整损坏矩阵。
\end{enumerate}

\section{可复现性与产物索引}

本轮累计 GPU 训练时间为 6.73 小时，结果完整性检查为 270/270。所有 checkpoint、逐帧预测、逐 episode 指标与 run metadata 均保存在实验目录中。

\noindent\begin{tabularx}{\linewidth}{p{39mm}Y}
\toprule
内容 & 路径\\
\midrule
实验定义与总览 & \path{runs/quality_ablation_v1_2_0/study.json}\\
结构化汇总 & \path{runs/quality_ablation_v1_2_0/results.csv}\\
交互式 HTML 报告 & \path{runs/quality_ablation_v1_2_0/report.html}\\
单次运行记录 & \path{runs/quality_ablation_v1_2_0/<dataset>/<model>/<condition>/seed-*/}\\
\bottomrule
\end{tabularx}

\vfill
\begin{center}
\small\color{muted}
本报告由 Sync2Act v1.2.0 实验输出生成。数值以结构化结果文件为准。
\end{center}

\end{document}
